% --- Archon physics formalization source begin ---
% archon:physics
% archon:covers PhyXMiniProblems/problem_phyx_mini_0860.lean
% archon:source-report reports/phyx_mini/problem_phyx_mini_0860.source.json

\chapter{Physics problem phyx\_mini\_0860}
\label{ch:PhyXMiniProblems_problem_phyx_mini_0860}

\paragraph{Problem source.}
\#\# Physical scenario

The electric potential at the dot in figure is 3140 V.

\#\# Question

What is charge q?

\#\# Answer choices

- A: A: 12nC

- B: B: 14nC

- C: C: 20nC

- D: D: 10nC

\#\# Figure caption (auxiliary; use the image as primary evidence)

The image depicts a physics diagram involving point charges and distances:

1. **Objects/Charges:**
   - There are three charged points, represented as circles with plus or minus signs inside.
   - The circle at the top left is labeled "q" and has a positive charge.
   - Below it, along the vertical dashed line, there is another circle labeled "5.0 nC" with a positive charge.
   - To the right of the top circle, there is a circle labeled "-5.0 nC" with a negative charge.

2. **Distances:**
   - The vertical distance between the two positive charges is labeled "4.0 cm."
   - The horizontal distance between the top positive charge and the negative charge is labeled "2.0 cm."

3. **Additional Elements:**
   - There is a small black dot located to the right of the vertical line near the bottom-right side of the diagram.

4. **Lines and Orientation:**
   - Dashed lines are used to represent the distances between the charges.

This diagram is likely used to illustrate electric forces or field calculations involving these point charges.

\#\# Dataset metadata

Electrostatics; Physical Model Grounding Reasoning, Multi-Formula Reasoning

\paragraph{Recorded answer/context.}
D: D: 10nC

\paragraph{Figure/image path.}
/root/proposal\_for\_physic/science-mango/phyx\_mini\_run/phyx\_data/test\_image/860.png

\paragraph{Formalization target.}
create a compiling Lean file with sorry bodies at `PhyXMiniProblems/problem\_phyx\_mini\_0860.lean`. The Lean declarations must preserve the physical quantities, dimensions or dimensional roles, figure labels, governing-law hypotheses, and final relation expressed by this problem.
Use Mathlib/Physlib names found through LeanExplore where available. If a physics API is missing, introduce faithful local abstractions rather than scalar placeholder aliases.

\begin{theorem}[Physics formalization target]
\label{thm:physics:phyx_mini_0860:target}
\lean{PhyXMiniProblems.ProblemPhyXMini0860.problem_phyx_mini_0860}
\uses{def:physics:phyx-mini-0860:phyxminiproblems-problemphyxmini0860-chargeinnanocoulombs, def:physics:phyx-mini-0860:phyxminiproblems-problemphyxmini0860-rectangularpointchargepotentialsetup, def:physics:phyx-mini-0860:phyxminiproblems-problemphyxmini0860-answerchoice, def:physics:phyx-mini-0860:phyxminiproblems-problemphyxmini0860-recordeddatasetanswer, def:physics:phyx-mini-0860:phyxminiproblems-problemphyxmini0860-matchessuppliedpotentialfigure, def:physics:phyx-mini-0860:phyxminiproblems-problemphyxmini0860-unknownchargeisoneofdisplayedchoices, def:physics:phyx-mini-0860:phyxminiproblems-problemphyxmini0860-usesintroductorycoulombconstant, def:physics:phyx-mini-0860:phyxminiproblems-problemphyxmini0860-hasphysicalpointchargeparameters, def:physics:phyx-mini-0860:phyxminiproblems-problemphyxmini0860-satisfiespointchargepotentiallaw, def:physics:phyx-mini-0860:phyxminiproblems-problemphyxmini0860-satisfieselectricpotentialsuperposition}
The assigned autoformalize agent should translate this physics problem into Lean declarations in the covered file, with theorem and lemma proof bodies written as `by sorry`.
\end{theorem}
\begin{proof}
This is an autoformalization task, not a proof task. Produce statements that can later be proved by the physics prover without weakening the physical model.
\end{proof}

% --- Archon declaration topology begin ---
\section{Declaration topology}
\begin{definition}[Length Quantity]
  \label{def:physics:phyx-mini-0860:phyxminiproblems-problemphyxmini0860-lengthquantity}
  \lean{PhyXMiniProblems.ProblemPhyXMini0860.LengthQuantity}
  A nonnegative, unit-independent physical length
\end{definition}
\begin{proof}
  This is the declared model definition of Length Quantity.
\end{proof}
\begin{definition}[Signed Charge Quantity]
  \label{def:physics:phyx-mini-0860:phyxminiproblems-problemphyxmini0860-signedchargequantity}
  \lean{PhyXMiniProblems.ProblemPhyXMini0860.SignedChargeQuantity}
  A signed, unit-independent physical electric charge
\end{definition}
\begin{proof}
  This is the declared model definition of Signed Charge Quantity.
\end{proof}
\begin{definition}[electric Potential Dimension]
  \label{def:physics:phyx-mini-0860:phyxminiproblems-problemphyxmini0860-electricpotentialdimension}
  \lean{PhyXMiniProblems.ProblemPhyXMini0860.electricPotentialDimension}
  The dimension M L² T⁻² C⁻¹ of electrostatic potential
\end{definition}
\begin{proof}
  This is the declared model definition of electric Potential Dimension.
\end{proof}
\begin{definition}[Electric Potential Quantity]
  \label{def:physics:phyx-mini-0860:phyxminiproblems-problemphyxmini0860-electricpotentialquantity}
  \lean{PhyXMiniProblems.ProblemPhyXMini0860.ElectricPotentialQuantity}
  \uses{def:physics:phyx-mini-0860:phyxminiproblems-problemphyxmini0860-electricpotentialdimension}
  A signed, unit-independent electrostatic potential
\end{definition}
\begin{proof}
  This is the declared model definition of Electric Potential Quantity.
\end{proof}
\begin{definition}[length Readout]
  \label{def:physics:phyx-mini-0860:phyxminiproblems-problemphyxmini0860-lengthreadout}
  \lean{PhyXMiniProblems.ProblemPhyXMini0860.lengthReadout}
  \uses{def:physics:phyx-mini-0860:phyxminiproblems-problemphyxmini0860-lengthquantity}
  Read a physical length in a selected Physlib length unit
\end{definition}
\begin{proof}
  This is the declared model definition of length Readout.
\end{proof}
\begin{definition}[length In Meters]
  \label{def:physics:phyx-mini-0860:phyxminiproblems-problemphyxmini0860-lengthinmeters}
  \lean{PhyXMiniProblems.ProblemPhyXMini0860.lengthInMeters}
  \uses{def:physics:phyx-mini-0860:phyxminiproblems-problemphyxmini0860-lengthquantity, def:physics:phyx-mini-0860:phyxminiproblems-problemphyxmini0860-lengthreadout}
  Coherent-SI metre readout of a physical length
\end{definition}
\begin{proof}
  This is the declared model definition of length In Meters.
\end{proof}
\begin{definition}[length In Centimeters]
  \label{def:physics:phyx-mini-0860:phyxminiproblems-problemphyxmini0860-lengthincentimeters}
  \lean{PhyXMiniProblems.ProblemPhyXMini0860.lengthInCentimeters}
  \uses{def:physics:phyx-mini-0860:phyxminiproblems-problemphyxmini0860-lengthquantity, def:physics:phyx-mini-0860:phyxminiproblems-problemphyxmini0860-lengthreadout}
  Centimetre readout used by the two dimension labels in the image
\end{definition}
\begin{proof}
  This is the declared model definition of length In Centimeters.
\end{proof}
\begin{definition}[charge Readout]
  \label{def:physics:phyx-mini-0860:phyxminiproblems-problemphyxmini0860-chargereadout}
  \lean{PhyXMiniProblems.ProblemPhyXMini0860.chargeReadout}
  \uses{def:physics:phyx-mini-0860:phyxminiproblems-problemphyxmini0860-signedchargequantity}
  Read a signed physical charge in a selected Physlib charge unit
\end{definition}
\begin{proof}
  This is the declared model definition of charge Readout.
\end{proof}
\begin{definition}[nanocoulomb Unit]
  \label{def:physics:phyx-mini-0860:phyxminiproblems-problemphyxmini0860-nanocoulombunit}
  \lean{PhyXMiniProblems.ProblemPhyXMini0860.nanocoulombUnit}
  The charge unit 10⁻⁹ C used by the three charge labels and choices
\end{definition}
\begin{proof}
  This is the declared model definition of nanocoulomb Unit.
\end{proof}
\begin{definition}[charge In Coulombs]
  \label{def:physics:phyx-mini-0860:phyxminiproblems-problemphyxmini0860-chargeincoulombs}
  \lean{PhyXMiniProblems.ProblemPhyXMini0860.chargeInCoulombs}
  \uses{def:physics:phyx-mini-0860:phyxminiproblems-problemphyxmini0860-signedchargequantity, def:physics:phyx-mini-0860:phyxminiproblems-problemphyxmini0860-chargereadout}
  Coherent-SI coulomb readout of a signed physical charge
\end{definition}
\begin{proof}
  This is the declared model definition of charge In Coulombs.
\end{proof}
\begin{definition}[charge In Nanocoulombs]
  \label{def:physics:phyx-mini-0860:phyxminiproblems-problemphyxmini0860-chargeinnanocoulombs}
  \lean{PhyXMiniProblems.ProblemPhyXMini0860.chargeInNanocoulombs}
  \uses{def:physics:phyx-mini-0860:phyxminiproblems-problemphyxmini0860-signedchargequantity, def:physics:phyx-mini-0860:phyxminiproblems-problemphyxmini0860-chargereadout, def:physics:phyx-mini-0860:phyxminiproblems-problemphyxmini0860-nanocoulombunit}
  Nanocoulomb readout used by the printed source labels and answer choices
\end{definition}
\begin{proof}
  This is the declared model definition of charge In Nanocoulombs.
\end{proof}
\begin{definition}[potential In Volts]
  \label{def:physics:phyx-mini-0860:phyxminiproblems-problemphyxmini0860-potentialinvolts}
  \lean{PhyXMiniProblems.ProblemPhyXMini0860.potentialInVolts}
  \uses{def:physics:phyx-mini-0860:phyxminiproblems-problemphyxmini0860-electricpotentialquantity}
  Coherent-SI readout of an electrostatic potential in volts
\end{definition}
\begin{proof}
  This is the declared model definition of potential In Volts.
\end{proof}
\begin{definition}[Figure Point]
  \label{def:physics:phyx-mini-0860:phyxminiproblems-problemphyxmini0860-figurepoint}
  \lean{PhyXMiniProblems.ProblemPhyXMini0860.FigurePoint}
  The four geometrical roles visible in the supplied rectangular diagram
\end{definition}
\begin{proof}
  This is the declared model definition of Figure Point.
\end{proof}
\begin{definition}[Source Charge]
  \label{def:physics:phyx-mini-0860:phyxminiproblems-problemphyxmini0860-sourcecharge}
  \lean{PhyXMiniProblems.ProblemPhyXMini0860.SourceCharge}
  The three point charges, named by their role in the image
\end{definition}
\begin{proof}
  This is the declared model definition of Source Charge.
\end{proof}
\begin{definition}[expected Source Point]
  \label{def:physics:phyx-mini-0860:phyxminiproblems-problemphyxmini0860-expectedsourcepoint}
  \lean{PhyXMiniProblems.ProblemPhyXMini0860.expectedSourcePoint}
  \uses{def:physics:phyx-mini-0860:phyxminiproblems-problemphyxmini0860-figurepoint, def:physics:phyx-mini-0860:phyxminiproblems-problemphyxmini0860-sourcecharge}
  The corner occupied by each of the three source charges
\end{definition}
\begin{proof}
  This is the declared model definition of expected Source Point.
\end{proof}
\begin{definition}[Figure Charge Sign]
  \label{def:physics:phyx-mini-0860:phyxminiproblems-problemphyxmini0860-figurechargesign}
  \lean{PhyXMiniProblems.ProblemPhyXMini0860.FigureChargeSign}
  The plus or minus glyph drawn inside a source circle
\end{definition}
\begin{proof}
  This is the declared model definition of Figure Charge Sign.
\end{proof}
\begin{definition}[expected Source Sign]
  \label{def:physics:phyx-mini-0860:phyxminiproblems-problemphyxmini0860-expectedsourcesign}
  \lean{PhyXMiniProblems.ProblemPhyXMini0860.expectedSourceSign}
  \uses{def:physics:phyx-mini-0860:phyxminiproblems-problemphyxmini0860-sourcecharge, def:physics:phyx-mini-0860:phyxminiproblems-problemphyxmini0860-figurechargesign}
  The sign glyph associated with each source in the primary image
\end{definition}
\begin{proof}
  This is the declared model definition of expected Source Sign.
\end{proof}
\begin{definition}[expected Printed Charge Text]
  \label{def:physics:phyx-mini-0860:phyxminiproblems-problemphyxmini0860-expectedprintedchargetext}
  \lean{PhyXMiniProblems.ProblemPhyXMini0860.expectedPrintedChargeText}
  \uses{def:physics:phyx-mini-0860:phyxminiproblems-problemphyxmini0860-sourcecharge}
  Literal text printed next to each source circle
\end{definition}
\begin{proof}
  This is the declared model definition of expected Printed Charge Text.
\end{proof}
\begin{definition}[Rectangular Potential Figure]
  \label{def:physics:phyx-mini-0860:phyxminiproblems-problemphyxmini0860-rectangularpotentialfigure}
  \lean{PhyXMiniProblems.ProblemPhyXMini0860.RectangularPotentialFigure}
  \uses{def:physics:phyx-mini-0860:phyxminiproblems-problemphyxmini0860-figurepoint, def:physics:phyx-mini-0860:phyxminiproblems-problemphyxmini0860-sourcecharge, def:physics:phyx-mini-0860:phyxminiproblems-problemphyxmini0860-figurechargesign}
  Literal presentation data transcribed from image 860.png
\end{definition}
\begin{proof}
  This is the declared model definition of Rectangular Potential Figure.
\end{proof}
\begin{definition}[Rectangular Point Charge Potential Setup]
  \label{def:physics:phyx-mini-0860:phyxminiproblems-problemphyxmini0860-rectangularpointchargepotentialsetup}
  \lean{PhyXMiniProblems.ProblemPhyXMini0860.RectangularPointChargePotentialSetup}
  \uses{def:physics:phyx-mini-0860:phyxminiproblems-problemphyxmini0860-lengthquantity, def:physics:phyx-mini-0860:phyxminiproblems-problemphyxmini0860-signedchargequantity, def:physics:phyx-mini-0860:phyxminiproblems-problemphyxmini0860-electricpotentialquantity, def:physics:phyx-mini-0860:phyxminiproblems-problemphyxmini0860-figurepoint, def:physics:phyx-mini-0860:phyxminiproblems-problemphyxmini0860-sourcecharge, def:physics:phyx-mini-0860:phyxminiproblems-problemphyxmini0860-rectangularpotentialfigure}
  The independent physical quantities and potential fields in the problem. charge .unknownQ is an unknown dimensionful charge, and is not defined from the recorded answer or from a displayed choice
\end{definition}
\begin{proof}
  This is the declared model definition of Rectangular Point Charge Potential Setup.
\end{proof}
\begin{definition}[source Position]
  \label{def:physics:phyx-mini-0860:phyxminiproblems-problemphyxmini0860-sourceposition}
  \lean{PhyXMiniProblems.ProblemPhyXMini0860.sourcePosition}
  \uses{def:physics:phyx-mini-0860:phyxminiproblems-problemphyxmini0860-sourcecharge, def:physics:phyx-mini-0860:phyxminiproblems-problemphyxmini0860-rectangularpointchargepotentialsetup}
  Position in the charge plane of a named source
\end{definition}
\begin{proof}
  This is the declared model definition of source Position.
\end{proof}
\begin{definition}[observation Position]
  \label{def:physics:phyx-mini-0860:phyxminiproblems-problemphyxmini0860-observationposition}
  \lean{PhyXMiniProblems.ProblemPhyXMini0860.observationPosition}
  \uses{def:physics:phyx-mini-0860:phyxminiproblems-problemphyxmini0860-rectangularpointchargepotentialsetup}
  Position marked by the black observation dot
\end{definition}
\begin{proof}
  This is the declared model definition of observation Position.
\end{proof}
\begin{definition}[position Vector In Meters]
  \label{def:physics:phyx-mini-0860:phyxminiproblems-problemphyxmini0860-positionvectorinmeters}
  \lean{PhyXMiniProblems.ProblemPhyXMini0860.positionVectorInMeters}
  Explicitly metre-valued coordinate vector associated with a planar point
\end{definition}
\begin{proof}
  This is the declared model definition of position Vector In Meters.
\end{proof}
\begin{definition}[displacement From Source In Meters]
  \label{def:physics:phyx-mini-0860:phyxminiproblems-problemphyxmini0860-displacementfromsourceinmeters}
  \lean{PhyXMiniProblems.ProblemPhyXMini0860.displacementFromSourceInMeters}
  \uses{def:physics:phyx-mini-0860:phyxminiproblems-problemphyxmini0860-sourcecharge, def:physics:phyx-mini-0860:phyxminiproblems-problemphyxmini0860-rectangularpointchargepotentialsetup, def:physics:phyx-mini-0860:phyxminiproblems-problemphyxmini0860-sourceposition, def:physics:phyx-mini-0860:phyxminiproblems-problemphyxmini0860-positionvectorinmeters}
  Displacement from a named source to an arbitrary planar point, in metres
\end{definition}
\begin{proof}
  This is the declared model definition of displacement From Source In Meters.
\end{proof}
\begin{definition}[source Distance In Meters]
  \label{def:physics:phyx-mini-0860:phyxminiproblems-problemphyxmini0860-sourcedistanceinmeters}
  \lean{PhyXMiniProblems.ProblemPhyXMini0860.sourceDistanceInMeters}
  \uses{def:physics:phyx-mini-0860:phyxminiproblems-problemphyxmini0860-sourcecharge, def:physics:phyx-mini-0860:phyxminiproblems-problemphyxmini0860-rectangularpointchargepotentialsetup, def:physics:phyx-mini-0860:phyxminiproblems-problemphyxmini0860-displacementfromsourceinmeters}
  Euclidean source-to-point separation in metres
\end{definition}
\begin{proof}
  This is the declared model definition of source Distance In Meters.
\end{proof}
\begin{definition}[Answer Choice]
  \label{def:physics:phyx-mini-0860:phyxminiproblems-problemphyxmini0860-answerchoice}
  \lean{PhyXMiniProblems.ProblemPhyXMini0860.AnswerChoice}
  Labels attached to the four displayed charge choices
\end{definition}
\begin{proof}
  This is the declared model definition of Answer Choice.
\end{proof}
\begin{definition}[charge In Nanocoulombs]
  \label{def:physics:phyx-mini-0860:phyxminiproblems-problemphyxmini0860-answerchoice-chargeinnanocoulombs}
  \lean{PhyXMiniProblems.ProblemPhyXMini0860.AnswerChoice.chargeInNanocoulombs}
  \uses{def:physics:phyx-mini-0860:phyxminiproblems-problemphyxmini0860-chargeinnanocoulombs, def:physics:phyx-mini-0860:phyxminiproblems-problemphyxmini0860-answerchoice}
  Nanocoulomb value printed by each answer choice
\end{definition}
\begin{proof}
  This is the declared model definition of charge In Nanocoulombs.
\end{proof}
\begin{definition}[recorded Dataset Answer]
  \label{def:physics:phyx-mini-0860:phyxminiproblems-problemphyxmini0860-recordeddatasetanswer}
  \lean{PhyXMiniProblems.ProblemPhyXMini0860.recordedDatasetAnswer}
  \uses{def:physics:phyx-mini-0860:phyxminiproblems-problemphyxmini0860-answerchoice}
  Answer label recorded in the supplied dataset metadata
\end{definition}
\begin{proof}
  This is the declared model definition of recorded Dataset Answer.
\end{proof}
\begin{definition}[Rounds To Nearest Ten Volts]
  \label{def:physics:phyx-mini-0860:phyxminiproblems-problemphyxmini0860-roundstonearesttenvolts}
  \lean{PhyXMiniProblems.ProblemPhyXMini0860.RoundsToNearestTenVolts}
  Agreement with a voltage printed to the nearest ten volts
\end{definition}
\begin{proof}
  This is the declared model definition of Rounds To Nearest Ten Volts.
\end{proof}
\begin{definition}[Matches Supplied Potential Figure]
  \label{def:physics:phyx-mini-0860:phyxminiproblems-problemphyxmini0860-matchessuppliedpotentialfigure}
  \lean{PhyXMiniProblems.ProblemPhyXMini0860.MatchesSuppliedPotentialFigure}
  \uses{def:physics:phyx-mini-0860:phyxminiproblems-problemphyxmini0860-lengthinmeters, def:physics:phyx-mini-0860:phyxminiproblems-problemphyxmini0860-lengthincentimeters, def:physics:phyx-mini-0860:phyxminiproblems-problemphyxmini0860-chargeincoulombs, def:physics:phyx-mini-0860:phyxminiproblems-problemphyxmini0860-chargeinnanocoulombs, def:physics:phyx-mini-0860:phyxminiproblems-problemphyxmini0860-potentialinvolts, def:physics:phyx-mini-0860:phyxminiproblems-problemphyxmini0860-expectedsourcepoint, def:physics:phyx-mini-0860:phyxminiproblems-problemphyxmini0860-expectedsourcesign, def:physics:phyx-mini-0860:phyxminiproblems-problemphyxmini0860-expectedprintedchargetext, def:physics:phyx-mini-0860:phyxminiproblems-problemphyxmini0860-rectangularpointchargepotentialsetup, def:physics:phyx-mini-0860:phyxminiproblems-problemphyxmini0860-observationposition, def:physics:phyx-mini-0860:phyxminiproblems-problemphyxmini0860-positionvectorinmeters, def:physics:phyx-mini-0860:phyxminiproblems-problemphyxmini0860-roundstonearesttenvolts}
  The literal figure features, their dimensional readouts, the reported potential, and their association with the physical sources. This contains no numerical value for the unknown charge
\end{definition}
\begin{proof}
  This is the declared model definition of Matches Supplied Potential Figure.
\end{proof}
\begin{definition}[Unknown Charge Is One Of Displayed Choices]
  \label{def:physics:phyx-mini-0860:phyxminiproblems-problemphyxmini0860-unknownchargeisoneofdisplayedchoices}
  \lean{PhyXMiniProblems.ProblemPhyXMini0860.UnknownChargeIsOneOfDisplayedChoices}
  \uses{def:physics:phyx-mini-0860:phyxminiproblems-problemphyxmini0860-chargeinnanocoulombs, def:physics:phyx-mini-0860:phyxminiproblems-problemphyxmini0860-rectangularpointchargepotentialsetup, def:physics:phyx-mini-0860:phyxminiproblems-problemphyxmini0860-answerchoice}
  The multiple-choice convention restricts the unknown to one of all four displayed values. It does not select D or assert the requested 10 nC conclusion
\end{definition}
\begin{proof}
  This is the declared model definition of Unknown Charge Is One Of Displayed Choices.
\end{proof}
\begin{definition}[Uses Introductory Coulomb Constant]
  \label{def:physics:phyx-mini-0860:phyxminiproblems-problemphyxmini0860-usesintroductorycoulombconstant}
  \lean{PhyXMiniProblems.ProblemPhyXMini0860.UsesIntroductoryCoulombConstant}
  \uses{def:physics:phyx-mini-0860:phyxminiproblems-problemphyxmini0860-rectangularpointchargepotentialsetup}
  The school-physics approximation to Coulomb's constant used with the displayed three-significant-figure potential. The scalar is the coherent-SI readout in N m²/C²
\end{definition}
\begin{proof}
  This is the declared model definition of Uses Introductory Coulomb Constant.
\end{proof}
\begin{definition}[Has Physical Point Charge Parameters]
  \label{def:physics:phyx-mini-0860:phyxminiproblems-problemphyxmini0860-hasphysicalpointchargeparameters}
  \lean{PhyXMiniProblems.ProblemPhyXMini0860.HasPhysicalPointChargeParameters}
  \uses{def:physics:phyx-mini-0860:phyxminiproblems-problemphyxmini0860-lengthinmeters, def:physics:phyx-mini-0860:phyxminiproblems-problemphyxmini0860-rectangularpointchargepotentialsetup, def:physics:phyx-mini-0860:phyxminiproblems-problemphyxmini0860-observationposition, def:physics:phyx-mini-0860:phyxminiproblems-problemphyxmini0860-sourcedistanceinmeters}
  Positivity and separation conditions selecting the physical branch
\end{definition}
\begin{proof}
  This is the declared model definition of Has Physical Point Charge Parameters.
\end{proof}
\begin{definition}[Satisfies Point Charge Potential Law]
  \label{def:physics:phyx-mini-0860:phyxminiproblems-problemphyxmini0860-satisfiespointchargepotentiallaw}
  \lean{PhyXMiniProblems.ProblemPhyXMini0860.SatisfiesPointChargePotentialLaw}
  \uses{def:physics:phyx-mini-0860:phyxminiproblems-problemphyxmini0860-chargeincoulombs, def:physics:phyx-mini-0860:phyxminiproblems-problemphyxmini0860-potentialinvolts, def:physics:phyx-mini-0860:phyxminiproblems-problemphyxmini0860-rectangularpointchargepotentialsetup, def:physics:phyx-mini-0860:phyxminiproblems-problemphyxmini0860-sourceposition, def:physics:phyx-mini-0860:phyxminiproblems-problemphyxmini0860-sourcedistanceinmeters}
  With electric potential zero at infinity, a point charge Q at distance r contributes k Q / r. This all-points law is stated away from each source and contains no requested value of q
\end{definition}
\begin{proof}
  This is the declared model definition of Satisfies Point Charge Potential Law.
\end{proof}
\begin{definition}[Satisfies Electric Potential Superposition]
  \label{def:physics:phyx-mini-0860:phyxminiproblems-problemphyxmini0860-satisfieselectricpotentialsuperposition}
  \lean{PhyXMiniProblems.ProblemPhyXMini0860.SatisfiesElectricPotentialSuperposition}
  \uses{def:physics:phyx-mini-0860:phyxminiproblems-problemphyxmini0860-potentialinvolts, def:physics:phyx-mini-0860:phyxminiproblems-problemphyxmini0860-sourcecharge, def:physics:phyx-mini-0860:phyxminiproblems-problemphyxmini0860-rectangularpointchargepotentialsetup}
  The total scalar potential is the sum of all three source potentials
\end{definition}
\begin{proof}
  This is the declared model definition of Satisfies Electric Potential Superposition.
\end{proof}
\begin{lemma}[source Distances At Observation]
  \label{lem:physics:phyx-mini-0860:phyxminiproblems-problemphyxmini0860-sourcedistancesatobservation}
  \lean{PhyXMiniProblems.ProblemPhyXMini0860.sourceDistancesAtObservation}
  \uses{def:physics:phyx-mini-0860:phyxminiproblems-problemphyxmini0860-rectangularpointchargepotentialsetup, def:physics:phyx-mini-0860:phyxminiproblems-problemphyxmini0860-observationposition, def:physics:phyx-mini-0860:phyxminiproblems-problemphyxmini0860-sourcedistanceinmeters, def:physics:phyx-mini-0860:phyxminiproblems-problemphyxmini0860-matchessuppliedpotentialfigure}
  The rectangle coordinates imply source-to-dot distances 2√5 cm, 4 cm, and 2 cm for the unknown, negative, and positive sources respectively
\end{lemma}
\begin{proof}
  The conclusion follows from the governing relations and figure readouts named in the statement of source Distances At Observation.
\end{proof}
\begin{lemma}[total Potential At Observation eq coulomb Sum]
  \label{lem:physics:phyx-mini-0860:phyxminiproblems-problemphyxmini0860-totalpotentialatobservation-eq-coulombsum}
  \lean{PhyXMiniProblems.ProblemPhyXMini0860.totalPotentialAtObservation_eq_coulombSum}
  \uses{def:physics:phyx-mini-0860:phyxminiproblems-problemphyxmini0860-chargeincoulombs, def:physics:phyx-mini-0860:phyxminiproblems-problemphyxmini0860-potentialinvolts, def:physics:phyx-mini-0860:phyxminiproblems-problemphyxmini0860-sourcecharge, def:physics:phyx-mini-0860:phyxminiproblems-problemphyxmini0860-rectangularpointchargepotentialsetup, def:physics:phyx-mini-0860:phyxminiproblems-problemphyxmini0860-observationposition, def:physics:phyx-mini-0860:phyxminiproblems-problemphyxmini0860-sourcedistanceinmeters, def:physics:phyx-mini-0860:phyxminiproblems-problemphyxmini0860-hasphysicalpointchargeparameters, def:physics:phyx-mini-0860:phyxminiproblems-problemphyxmini0860-satisfiespointchargepotentiallaw, def:physics:phyx-mini-0860:phyxminiproblems-problemphyxmini0860-satisfieselectricpotentialsuperposition}
  Point-charge potential and superposition give the source-by-source expression for the potential at the observation dot. Its right-hand side contains only independent source data, geometry, and the governing Coulomb constant
\end{lemma}
\begin{proof}
  The conclusion follows from the governing relations and figure readouts named in the statement of total Potential At Observation eq coulomb Sum.
\end{proof}
% --- Archon declaration topology end ---
% --- Archon physics formalization source end ---
